% Appendix A — 7-Dimension BIO Label Schema

\chapter{Seven-Dimension BIO Label Schema and Clinical Examples}
\label{AppendixA}

This appendix provides the complete specification of the 15-class BIO token label schema used in Kent-AI, along with annotated clinical examples illustrating each label class.

%----------------------------------------------------------------------------------------

\section{BIO Label Specification}

\begin{table}[h]
\caption[Complete BIO Label Schema]{Complete 15-Class BIO Token Label Schema with Clinical Definitions.}
\label{tab:fullschema}
\centering
\begin{tabular}{lllp{6cm}}
\toprule
\tabhead{Label} & \tabhead{Dimension} & \tabhead{Abbrev.} & \tabhead{Clinical Definition} \\
\midrule
B-MENT & Mental/Emotional & MENT & Beginning of a mental or emotional symptom span. Includes anxiety, fears, grief, irritability, restlessness, sadness, despair, and all mental/emotional states. \\
I-MENT & Mental/Emotional & MENT & Continuation of a MENT span. \\
\midrule
B-LOC & Location & LOC & Beginning of an anatomical location span. Includes organ names, body regions, laterality (left, right, bilateral), and directional descriptors (anterior, posterior). \\
I-LOC & Location & LOC & Continuation of a LOC span. \\
\midrule
B-SEN & Sensation & SEN & Beginning of a sensation or pain quality span. Includes descriptors such as throbbing, burning, pressing, stabbing, stitching, cramping, numb. \\
I-SEN & Sensation & SEN & Continuation of a SEN span. \\
\midrule
B-MOD\_AGG & Aggravation & MOD\_AGG & Beginning of an aggravation modality span. Describes conditions or actions that make symptoms worse (e.g., cold air, motion, evening, being alone). \\
I-MOD\_AGG & Aggravation & MOD\_AGG & Continuation of a MOD\_AGG span. \\
\midrule
B-MOD\_AMEL & Amelioration & MOD\_AMEL & Beginning of an amelioration modality span. Describes conditions that relieve symptoms (e.g., warmth, rest, company, pressure). \\
I-MOD\_AMEL & Amelioration & MOD\_AMEL & Continuation of a MOD\_AMEL span. \\
\midrule
B-CONC & Concomitant & CONC & Beginning of a concomitant symptom span. Describes symptoms that consistently accompany the chief complaint but are anatomically or etiologically distinct. \\
I-CONC & Concomitant & CONC & Continuation of a CONC span. \\
\midrule
B-TEMP & Temporal & TEMP & Beginning of a temporal or periodicity span. Includes time of day, day of week, season, duration, frequency, and onset descriptions. \\
I-TEMP & Temporal & TEMP & Continuation of a TEMP span. \\
\midrule
O & Outside & O & Token is not part of any symptom entity. \\
\bottomrule
\end{tabular}
\end{table}

%----------------------------------------------------------------------------------------

\section{Annotated Clinical Examples}

\subsection{Example 1: Mental/Emotional + Aggravation + Temporal}

\textbf{Patient narrative}: \textit{``I feel terribly anxious and restless, especially in the evening when I am alone.''}

\begin{verbatim}
Token          Label
-----------    ----------
I              O
feel           O
terribly       O
anxious        B-MENT
and            I-MENT
restless       I-MENT
,              O
especially     O
in             B-TEMP
the            I-TEMP
evening        I-TEMP
when           B-MOD_AGG
I              I-MOD_AGG
am             I-MOD_AGG
alone          I-MOD_AGG
.              O
\end{verbatim}

\textbf{Matched Rubric}: \texttt{MIND > ANXIETY > evening > alone, when}

\subsection{Example 2: Location + Sensation + Amelioration}

\textbf{Patient narrative}: \textit{``There is a severe throbbing pain in my right temple that is better from firm pressure.''}

\begin{verbatim}
Token          Label
-----------    ----------
There          O
is             O
a              O
severe         B-SEN
throbbing      I-SEN
pain           I-SEN
in             O
my             O
right          B-LOC
temple         I-LOC
that           O
is             O
better         B-MOD_AMEL
from           I-MOD_AMEL
firm           I-MOD_AMEL
pressure       I-MOD_AMEL
.              O
\end{verbatim}

\textbf{Matched Rubric}: \texttt{HEAD > PAIN > temples > right > pressing > amel.}

\subsection{Example 3: Concomitant + Temporal}

\textbf{Patient narrative}: \textit{``My headaches come every afternoon around 3 to 4 PM and are always accompanied by nausea.''}

\begin{verbatim}
Token          Label
-----------    ----------
My             O
headaches      O
come           O
every          B-TEMP
afternoon      I-TEMP
around         I-TEMP
3              I-TEMP
to             I-TEMP
4              I-TEMP
PM             I-TEMP
and            O
are            O
always         O
accompanied    O
by             O
nausea         B-CONC
.              O
\end{verbatim}
